\color{unidarkgreen}{\section[(обязательное) Уставки функций РЗА]{(обязательное)\\Уставки функций РЗА}\label{app:set}}
\color{black}

\setlength{\extrarowheight}{0.06cm}

\par
Уставки ДТЗ представлены в таблице \ref{app_set:dzt}.

\footnotesize
\begin{longtable}{|>{\centering\arraybackslash}m{5.3cm}|>{\centering\arraybackslash}m{3.3cm}|>{\centering\arraybackslash}m{4.2cm}|>{\centering\arraybackslash}m{1.8cm}|>{\centering\arraybackslash}m{1cm}|}
% Первая страница: полный заголовок
\caption{Параметры для настройки ДТЗ\hfill\vspace{-0.5\baselineskip}}\label{app_set:dzt}\\
\hline
\rowcolor{gray!20}
Параметр & Условное обозначение на схеме & Значение/ Диапазон & Единица измерения & Шаг \\
\hline
\endfirsthead

% Промежуточные страницы: «Продолжение таблицы»
\caption*{\hspace{3pt}\emph{Продолжение таблицы \ref{app_set:dzt}\hfill\vspace{-0.5\baselineskip}}} \\ % сделано по ГОСТ 2.105 п.6.8.7

\hline
\rowcolor{gray!20}
Параметр & Условное обозначение на схеме & Значение/ Диапазон & Единица измерения & Шаг \\
\hline
\endhead

%===m>ДТЗ|general+-
\multicolumn{5}{|c|}{ ДТЗ } \\ \hline 
\raggedright ДТЗ фВ & \centering ДТЗ фВ & \centering Выведено \\ Введено & \centering --- & \centering \arraybackslash --- \\
\hline
\raggedright Iср диф & \centering Iср диф & \centering 0,100 ... 2,000 & \centering о.е. & \centering \arraybackslash 0,001 \\
\hline
\raggedright \textcolor{red}{FullDescription НЕ НАЙДЕН!} & \centering Кв(Iср диф) & \centering 0,00 ... 1,00 & \centering о.е. & \centering \arraybackslash 0,01 \\
\hline
\multicolumn{5}{|c|}{ КЦТ } \\ \hline 
\raggedright Iнеиспр ДТЗ & \centering Iнеиспр ДТЗ & \centering 0,04 ... 30,00 & \centering о.е. & \centering \arraybackslash 0,01 \\
\hline
\raggedright \textcolor{red}{FullDescription НЕ НАЙДЕН!} & \centering Кв (Iнеиспр) & \centering 0,00 ... 1,00 & \centering  & \centering \arraybackslash 0,01 \\
\hline
\raggedright Тнеиспр ДТЗ & \centering Тнеиспр ДТЗ & \centering 0 ... 300000 & \centering мс & \centering \arraybackslash 5 \\
\hline
\raggedright Контр изм цеп ДТЗ & \centering Контр изм цеп ДТЗ & \centering Введено \\ Выведено & \centering --- & \centering \arraybackslash --- \\
\hline
\raggedright \textcolor{red}{FullDescription НЕ НАЙДЕН!} & \centering фВ & \centering Введено \\ Выведено & \centering --- & \centering \arraybackslash --- \\
\hline
\multicolumn{5}{|c|}{ ДТЗт } \\ \hline 
\raggedright \textcolor{red}{FullDescription НЕ НАЙДЕН!} & \centering Iср диф & \centering 0,100 ... 2,000 & \centering о.е. & \centering \arraybackslash 0,001 \\
\hline
\raggedright Iт1 & \centering Iт1 & \centering 0,200 ... 1,500 & \centering о.е. & \centering \arraybackslash 0,001 \\
\hline
\raggedright Кт1 & \centering Кт1 & \centering 0,100 ... 1,000 & \centering о.е. & \centering \arraybackslash 0,001 \\
\hline
\raggedright Iт2 & \centering Iт2 & \centering 1,000 ... 3,000 & \centering о.е. & \centering \arraybackslash 0,001 \\
\hline
\raggedright Кт2 & \centering Кт2 & \centering 0,200 ... 1,000 & \centering о.е. & \centering \arraybackslash 0,001 \\
\hline
\raggedright \textcolor{red}{FullDescription НЕ НАЙДЕН!} & \centering Кв (Iср диф) & \centering 0,00 ... 1,00 & \centering  & \centering \arraybackslash 0,01 \\
\hline
\raggedright \textcolor{red}{FullDescription НЕ НАЙДЕН!} & \centering фВ & \centering Введено \\ Выведено & \centering --- & \centering \arraybackslash --- \\
\hline
\raggedright ДТЗт & \centering ДТЗт & \centering Введено \\ Выведено & \centering --- & \centering \arraybackslash --- \\
\hline
\multicolumn{5}{|c|}{ ДТО } \\ \hline 
\raggedright Iср ДТО & \centering Iср ДТО & \centering 1,000 ... 20,000 & \centering о.е. & \centering \arraybackslash 0,001 \\
\hline
\raggedright \textcolor{red}{FullDescription НЕ НАЙДЕН!} & \centering Кв (Iср) & \centering 0,00 ... 1,00 & \centering  & \centering \arraybackslash 0,01 \\
\hline
\raggedright \textcolor{red}{FullDescription НЕ НАЙДЕН!} & \centering фВ & \centering Введено \\ Выведено & \centering --- & \centering \arraybackslash --- \\
\hline
\raggedright ДТО & \centering ДТО & \centering Введено \\ Выведено & \centering --- & \centering \arraybackslash --- \\
\hline
\multicolumn{5}{|c|}{ Д2г } \\ \hline 
\raggedright I2г/I1г ДТЗт & \centering I2г/I1г ДТЗт & \centering 0,050 ... 1,000 & \centering о.е. & \centering \arraybackslash 0,001 \\
\hline
\raggedright \textcolor{red}{FullDescription НЕ НАЙДЕН!} & \centering Кв (I2(5)г/I1г) & \centering 0,00 ... 1,00 & \centering  & \centering \arraybackslash 0,01 \\
\hline
\raggedright \textcolor{red}{FullDescription НЕ НАЙДЕН!} & \centering Iср диф & \centering 0,000 ... 2,000 & \centering о.е. & \centering \arraybackslash 0,001 \\
\hline
\raggedright \textcolor{red}{FullDescription НЕ НАЙДЕН!} & \centering Кв (Iср диф) & \centering 0,00 ... 1,00 & \centering  & \centering \arraybackslash 0,01 \\
\hline
\raggedright \textcolor{red}{FullDescription НЕ НАЙДЕН!} & \centering 0.01 с & \centering 0 ... 10 & \centering мc & \centering \arraybackslash 5 \\
\hline
\raggedright \textcolor{red}{FullDescription НЕ НАЙДЕН!} & \centering 0.02 с & \centering 0 ... 20 & \centering мc & \centering \arraybackslash 5 \\
\hline
\raggedright \textcolor{red}{FullDescription НЕ НАЙДЕН!} & \centering фВ & \centering Введено \\ Выведено & \centering --- & \centering \arraybackslash --- \\
\hline
\raggedright Перекр блок 2г & \centering Перекр блок 2г & \centering Откл \\ 1 из 3 \\ 2 из 3 & \centering --- & \centering \arraybackslash --- \\
\hline
\raggedright Тперекр блок 2г & \centering Тперекр блок 2г & \centering 100 ... 1000 & \centering мc & \centering \arraybackslash 5 \\
\hline
\raggedright Блок 2г ДТЗт & \centering Блок 2г ДТЗт & \centering Введено \\ Выведено & \centering --- & \centering \arraybackslash --- \\
\hline
\multicolumn{5}{|c|}{ Д5г } \\ \hline 
\raggedright I5г/I1г ДТЗт & \centering I5г/I1г ДТЗт & \centering 0,050 ... 1,000 & \centering о.е. & \centering \arraybackslash 0,001 \\
\hline
\raggedright \textcolor{red}{FullDescription НЕ НАЙДЕН!} & \centering Кв (I2(5)г/I1г) & \centering 0,00 ... 1,00 & \centering  & \centering \arraybackslash 0,01 \\
\hline
\raggedright \textcolor{red}{FullDescription НЕ НАЙДЕН!} & \centering Iср диф & \centering 0,000 ... 2,000 & \centering о.е. & \centering \arraybackslash 0,001 \\
\hline
\raggedright \textcolor{red}{FullDescription НЕ НАЙДЕН!} & \centering Кв (Iср диф) & \centering 0,00 ... 1,00 & \centering  & \centering \arraybackslash 0,01 \\
\hline
\raggedright \textcolor{red}{FullDescription НЕ НАЙДЕН!} & \centering 0.01 с & \centering 0 ... 10 & \centering мc & \centering \arraybackslash 5 \\
\hline
\raggedright \textcolor{red}{FullDescription НЕ НАЙДЕН!} & \centering 0.02 с & \centering 0 ... 20 & \centering мc & \centering \arraybackslash 5 \\
\hline
\raggedright \textcolor{red}{FullDescription НЕ НАЙДЕН!} & \centering фВ & \centering Введено \\ Выведено & \centering --- & \centering \arraybackslash --- \\
\hline
\raggedright Перекр блок 5г & \centering Перекр блок 5г & \centering Откл \\ 1 из 3 \\ 2 из 3 & \centering --- & \centering \arraybackslash --- \\
\hline
\raggedright Тперекр блок 5г & \centering Тперекр блок 5г & \centering 100 ... 1000 & \centering мc & \centering \arraybackslash 5 \\
\hline
\raggedright Блок 5г ДТЗт & \centering Блок 5г ДТЗт & \centering Введено \\ Выведено & \centering --- & \centering \arraybackslash --- \\
\hline
%===m
\end{longtable}
\normalsize


%%%%%%%%%%%%%%%%%%%%%%%%%%%%%%%%%%%%%%% МТЗ НН

\par
Уставки МТЗ НН представлены в таблице \ref{app_set:mtz}.

\footnotesize
\begin{longtable}{|>{\centering\arraybackslash}m{5.3cm}|>{\centering\arraybackslash}m{3.3cm}|>{\centering\arraybackslash}m{4.2cm}|>{\centering\arraybackslash}m{1.8cm}|>{\centering\arraybackslash}m{1cm}|}
% Первая страница: полный заголовок
\caption{Параметры для настройки МТЗ НН\hfill\vspace{-0.5\baselineskip}}\label{app_set:mtz}\\
\hline
\rowcolor{gray!20}
Параметр & Условное обозначение на схеме & Значение/ Диапазон & Единица измерения & Шаг \\
\hline
\endfirsthead

% Промежуточные страницы: «Продолжение таблицы»
\caption*{\hspace{3pt}\emph{Продолжение таблицы \ref{app_set:mtz}\hfill\vspace{-0.5\baselineskip}}} \\ % сделано по ГОСТ 2.105 п.6.8.7

\hline
\rowcolor{gray!20}
Параметр & Условное обозначение на схеме & Значение/ Диапазон & Единица измерения & Шаг \\
\hline
\endhead

%===m>МТЗ НН|general+-
\multicolumn{5}{|c|}{ МТЗ-1 } \\ \hline 
\raggedright Iср МТЗ-1 НН & \centering Iср МТЗ-1 НН & \centering 0,04 ... 30,00 & \centering о.е. & \centering \arraybackslash 0,01 \\
\hline
\raggedright Тср МТЗ-1 НН & \centering Тср МТЗ-1 НН & \centering 0 ... 300000 & \centering мс & \centering \arraybackslash 5 \\
\hline
\raggedright МТЗ-1 НН & \centering МТЗ-1 НН & \centering Введено \\ Выведено & \centering --- & \centering \arraybackslash --- \\
\hline
\raggedright Режим КПОН МТЗ-1 НН & \centering Режим КПОН МТЗ-1 НН & \centering Введено \\ Выведено & \centering --- & \centering \arraybackslash --- \\
\hline
\raggedright Режим БНТ МТЗ-1 НН & \centering Режим БНТ МТЗ-1 НН & \centering Введено \\ Выведено & \centering --- & \centering \arraybackslash --- \\
\hline
\multicolumn{5}{|c|}{ МТЗ-2 } \\ \hline 
\raggedright Iср МТЗ-2 НН & \centering Iср МТЗ-2 НН & \centering 0,04 ... 30,00 & \centering о.е. & \centering \arraybackslash 0,01 \\
\hline
\raggedright Тср МТЗ-2 НН & \centering Тср МТЗ-2 НН & \centering 0 ... 300000 & \centering мс & \centering \arraybackslash 5 \\
\hline
\raggedright МТЗ-2 НН & \centering МТЗ-2 НН & \centering Введено \\ Выведено & \centering --- & \centering \arraybackslash --- \\
\hline
\raggedright Режим КПОН МТЗ-2 НН & \centering Режим КПОН МТЗ-2 НН & \centering Введено \\ Выведено & \centering --- & \centering \arraybackslash --- \\
\hline
\raggedright Режим БНТ МТЗ-2 НН & \centering Режим БНТ МТЗ-2 НН & \centering Введено \\ Выведено & \centering --- & \centering \arraybackslash --- \\
\hline
\multicolumn{5}{|c|}{ ТО } \\ \hline 
\raggedright Iср ТО НН & \centering Iср ТО НН & \centering 0,04 ... 30,00 & \centering о.е. & \centering \arraybackslash 0,01 \\
\hline
\raggedright Тср ТО НН & \centering Тср ТО НН & \centering 0 ... 300000 & \centering мс & \centering \arraybackslash 5 \\
\hline
\raggedright  ТО НН & \centering  ТО НН & \centering Введено \\ Выведено & \centering --- & \centering \arraybackslash --- \\
\hline
\raggedright Режим БНТ ТО НН & \centering Режим БНТ ТО НН & \centering Введено \\ Выведено & \centering --- & \centering \arraybackslash --- \\
\hline
\multicolumn{5}{|c|}{ БНТ } \\ \hline 
\raggedright I2г/I1г МТЗ НН & \centering I2г/I1г МТЗ НН & \centering 0,00 ... 1,00 & \centering о.е. & \centering \arraybackslash 0,01 \\
\hline
\raggedright \textcolor{red}{FullDescription НЕ НАЙДЕН!} & \centering kв=0.95 & \centering 0,00 ... 1,00 & \centering  & \centering \arraybackslash 0,01 \\
\hline
\raggedright 0.04Iном (Iразр БНТ МТЗ НН) & \centering 0.04Iном (Iразр БНТ МТЗ НН) & \centering 0,04 ... 30,00 & \centering о.е & \centering \arraybackslash 0,01 \\
\hline
\raggedright \textcolor{red}{FullDescription НЕ НАЙДЕН!} & \centering kв=1 & \centering 0,00 ... 1,00 & \centering  & \centering \arraybackslash 0,01 \\
\hline
\raggedright БНТ МТЗ НН & \centering БНТ МТЗ НН & \centering Введено \\ Выведено & \centering --- & \centering \arraybackslash --- \\
\hline
\raggedright Iблок БНТ МТЗ НН & \centering Iблок БНТ МТЗ НН & \centering 0,04 ... 30,00 & \centering о.е. & \centering \arraybackslash 0,01 \\
\hline
%===m
\end{longtable}
\normalsize


%%%%%%%%%%%%%%%%%%%%%%%%%%%%%%%%%%%%%%% КПОН

\par
Уставки КПОН представлены в таблице \ref{app_set:kpon}.

\footnotesize
\begin{longtable}{|>{\centering\arraybackslash}m{5.3cm}|>{\centering\arraybackslash}m{3.3cm}|>{\centering\arraybackslash}m{4.2cm}|>{\centering\arraybackslash}m{1.8cm}|>{\centering\arraybackslash}m{1cm}|}
% Первая страница: полный заголовок
\caption{Параметры для настройки КПОН\hfill\vspace{-0.5\baselineskip}}\label{app_set:kpon}\\
\hline
\rowcolor{gray!20}
Параметр & Условное обозначение на схеме & Значение/ Диапазон & Единица измерения & Шаг \\
\hline
\endfirsthead

% Промежуточные страницы: «Продолжение таблицы»
\caption*{\hspace{3pt}\emph{Продолжение таблицы \ref{app_set:kpon}\hfill\vspace{-0.5\baselineskip}}} \\ % сделано по ГОСТ 2.105 п.6.8.7

\hline
\rowcolor{gray!20}
Параметр & Условное обозначение на схеме & Значение/ Диапазон & Единица измерения & Шаг \\
\hline
\endhead

%===m>КПОН|general+-
\multicolumn{5}{|c|}{ КПОН НН1 } \\ \hline 
\raggedright U2 КПОН НН1 & \centering U2 КПОН НН1 & \centering 0,0050 ... 1,5000 & \centering о.е & \centering \arraybackslash 0,0001 \\
\hline
\raggedright Uл КПОН НН1 & \centering Uл КПОН НН1 & \centering 0,0050 ... 1,5000 & \centering о.е & \centering \arraybackslash 0,0001 \\
\hline
\raggedright Тип КПОН НН1 & \centering Тип КПОН НН1 & \centering Umin, U2 \\ Umin \\ Внеш & \centering --- & \centering \arraybackslash --- \\
\hline
\raggedright  КПОН НН1 & \centering  КПОН НН1 & \centering Введено \\ Выведено & \centering --- & \centering \arraybackslash --- \\
\hline
\raggedright Неиспр ТН КПОН НН1 & \centering Неиспр ТН КПОН НН1 & \centering Выведено \\ Пуск КПОН \\ Блок КПОН & \centering --- & \centering \arraybackslash --- \\
\hline
\multicolumn{5}{|c|}{ КПОН НН2 } \\ \hline 
\raggedright U2 КПОН НН2 & \centering U2 КПОН НН2 & \centering 0,0050 ... 1,5000 & \centering о.е & \centering \arraybackslash 0,0001 \\
\hline
\raggedright Uл КПОН НН2 & \centering Uл КПОН НН2 & \centering 0,0050 ... 1,5000 & \centering о.е & \centering \arraybackslash 0,0001 \\
\hline
\raggedright Тип КПОН НН2 & \centering Тип КПОН НН2 & \centering Umin, U2 \\ Umin \\ Внеш & \centering --- & \centering \arraybackslash --- \\
\hline
\raggedright  КПОН НН2 & \centering  КПОН НН2 & \centering Введено \\ Выведено & \centering --- & \centering \arraybackslash --- \\
\hline
\raggedright Неиспр ТН КПОН НН2 & \centering Неиспр ТН КПОН НН2 & \centering Выведено \\ Пуск КПОН \\ Блок КПОН & \centering --- & \centering \arraybackslash --- \\
\hline
\multicolumn{5}{|c|}{ КПОН НН } \\ \hline 
\raggedright U2 КПОН НН & \centering U2 КПОН НН & \centering 0,0050 ... 1,5000 & \centering о.е & \centering \arraybackslash 0,0001 \\
\hline
\raggedright Uл КПОН НН & \centering Uл КПОН НН & \centering 0,0050 ... 1,5000 & \centering о.е & \centering \arraybackslash 0,0001 \\
\hline
\raggedright Тип КПОН НН & \centering Тип КПОН НН & \centering Umin, U2 \\ Umin \\ Внеш & \centering --- & \centering \arraybackslash --- \\
\hline
\raggedright КПОН НН & \centering КПОН НН & \centering Введено \\ Выведено & \centering --- & \centering \arraybackslash --- \\
\hline
\raggedright Неиспр ТН КПОН НН & \centering Неиспр ТН КПОН НН & \centering Выведено \\ Пуск КПОН \\ Блок КПОН & \centering --- & \centering \arraybackslash --- \\
\hline
%===m
\end{longtable}
\normalsize